\section{Mô tả phần mềm cài đặt}

\subsection{Phạm vi trình bày}

Theo yêu cầu của báo cáo, phần này không mô tả cách cài đặt các ứng dụng phụ trợ như trình soạn mã hay công cụ xem cơ sở dữ liệu. Em cũng xin nói rõ rằng phần mềm này là sản phẩm thương mại đang được khai thác, nên em chỉ trình bày ở mức cách thức tổ chức và quy trình, không nêu địa chỉ máy chủ, thông số cấu hình, chuỗi kết nối hay bất kỳ chi tiết nào có thể ảnh hưởng tới hệ thống thật và tới dữ liệu của khách hàng.

\subsection{Môi trường và công cụ cần có}

Phần giao diện dùng Angular với TypeScript, trong đó màn hình lễ tân chạy trên Angular 9 còn màn hình quản lý viết lại trên Angular 18. Phần máy chủ dùng .NET và C\#, truy cập dữ liệu qua Entity Framework, lưu dữ liệu trên SQL Server, dùng Redis làm bộ đệm và khoá phân tán, và các dịch vụ mới được đóng gói bằng Docker. Việc dựng mã và chạy kiểm thử được tự động hoá trên một máy chủ tích hợp liên tục. Về kiểm thử, phần máy chủ dùng xUnit còn phần giao diện dùng Jasmine cùng Karma.

Điểm đáng lưu ý là hai phiên bản Angular cách nhau khá xa, phản ánh thực tế của một sản phẩm đã sống nhiều năm: phần cũ vẫn phải chạy và vẫn phải sửa, trong khi phần mới được dựng song song. Trong nhiều tuần em làm trên cả hai, và việc chuyển qua lại giữa hai lối viết của cùng một thư viện buộc em phân biệt được đâu là khái niệm cốt lõi của công cụ, đâu chỉ là cách viết của một phiên bản.

\subsection{Cách tổ chức mã nguồn và cách dựng phần mềm}

Cả hai phía đều được tổ chức theo tầng, và cách đặt tên thư mục nói lên thứ tự các tầng. Ở phía máy chủ, một dịch vụ chia thành các nhóm: nhóm dùng chung chứa tiện ích và hằng số, nhóm trình bày chứa đầu giao tiếp và mô tả dữ liệu vào ra, nhóm ứng dụng chứa quy tắc nghiệp vụ cùng phần xử lý lệnh và truy vấn, nhóm hạ tầng chứa phần nói chuyện với cơ sở dữ liệu, và một nhóm riêng cho kiểm thử. Cách đánh số thứ tự vào tên thư mục nhắc người đọc chiều phụ thuộc: nhóm bên ngoài được biết về nhóm bên trong, ngược lại thì không. Phía giao diện có cấu trúc tương tự. Khi làm tính năng khoá phòng, em phải thêm tệp vào gần như tất cả các thư mục này, và chính lần đó em hiểu vì sao mỗi tầng lại tồn tại.

Về cách dựng, phần giao diện cần cài các gói phụ thuộc theo bản khai báo đã khoá phiên bản, chọn tệp cấu hình môi trường, rồi chạy lệnh dựng để sinh bộ tệp tĩnh đem đặt lên máy chủ web; riêng màn hình lễ tân còn có bản đóng gói chạy như ứng dụng trên máy tính tại quầy, nên khi sửa những phần liên quan tới in ấn hay thiết bị ngoại vi thì phải thử trên cả hai dạng. Phần máy chủ thì phục hồi các gói phụ thuộc, dựng bộ giải pháp, chạy toàn bộ kiểm thử tự động, rồi xuất bản.

\subsection{Cách đưa thay đổi cơ sở dữ liệu}

Đây là phần em cẩn thận nhất, vì mã nguồn sai thì có thể đưa lại bản cũ, còn dữ liệu bị sửa sai thì không phải lúc nào cũng lấy lại được. Các thay đổi về cấu trúc dữ liệu như thêm bảng, thêm cột hay sửa câu lệnh sẵn được quản lý trong một dự án riêng, để mọi thay đổi đều có lịch sử và đều đi qua bước xem lại. Các thay đổi về nội dung dữ liệu được viết thành tập lệnh riêng, và với mỗi tập lệnh chuyển đổi em luôn viết kèm một tập lệnh đưa dữ liệu về trạng thái trước đó; em đã làm việc này cho cả lần tách quyền huỷ đơn và lần bổ sung dữ liệu về quốc gia, múi giờ, đơn vị tiền tệ. Thói quen này em học được từ nhóm chứ không phải tự nghĩ ra, và đến giờ em thấy đó là một trong những thói quen giá trị nhất mình mang về.

\subsection{Chuỗi môi trường và quy trình phát hành}

Một thay đổi đi qua năm chặng. Em làm và tự thử trên nhánh riêng của mình, đồng nghiệp xem lại mã, rồi mã được gộp vào môi trường phát triển chung nơi hệ thống tự dựng và tự chạy kiểm thử. Tiếp đó mã lên môi trường kiểm thử để bộ phận kiểm thử xác nhận, rồi vào nhánh chuẩn bị phát hành để chốt phạm vi và chuẩn bị các tập lệnh dữ liệu, và cuối cùng lên hệ thống thật theo lịch, kèm việc theo dõi sau phát hành.

Điều em muốn nhấn mạnh là chặng thứ ba có một đường quay lại: nếu bộ phận kiểm thử còn tìm ra lỗi thì phiếu về lại nhánh riêng, và trong sáu tháng khá nhiều phiếu của em đã đi theo đường đó ít nhất một lần. Ban đầu mỗi lần bị trả lại em thấy khá nản, về sau em hiểu rằng một phiếu bị trả lại ở môi trường kiểm thử thì rẻ hơn rất nhiều so với cùng lỗi đó bị phát hiện bởi khách hàng.

Hai điểm khác trong quy trình em thấy đáng ghi lại. Thứ nhất, tính năng mới không được bật cho tất cả khách hàng cùng lúc mà bật dần theo từng nhóm, và việc chuyển khách hàng từ bản quản lý cũ sang bản mới cũng làm theo từng bước có mốc thời gian rõ ràng, em tham gia phần này ở giai đoạn cuối kỳ thực tập. Thứ hai, khi phát hành phiên bản mới của phần giao diện thì trình duyệt của người dùng vẫn giữ bản cũ trong bộ đệm nên họ mở lên và không thấy thay đổi. Đây là loại vấn đề em hoàn toàn không nghĩ tới trước khi đi thực tập, vì khi tự làm dự án ở nhà thì chỉ cần tự xoá bộ đệm là xong.

\subsection{Kiểm thử tự động trong quy trình}

Trong sáu tháng em viết khoảng một nghìn chín trăm dòng kiểm thử tự động, phân bố trên bảy tệp, gồm năm tệp cho phần máy chủ và hai tệp cho phần giao diện.

Ở phía máy chủ, các tệp kiểm thử tập trung vào những chỗ có logic thật: phần sinh lệnh cho tính năng khoá phòng, các quy tắc kiểm tra dữ liệu của tính năng đó, phần thông báo theo ngôn ngữ, và phần tính toán thống kê. Em đặt tên các trường hợp kiểm thử theo dạng mô tả tình huống và kết quả mong đợi, ví dụ bỏ hết phòng mà không chọn phòng mới thì phải bị chặn, nên danh sách kết quả chạy kiểm thử đọc lên gần như một bản mô tả hành vi.

Ở phía giao diện, em viết kiểm thử cho phần kiểm tra dữ liệu khi đặt phòng, vì đó là chỗ nhiều quy tắc chồng nhau: phòng đang bị khoá, phòng đang có khách chưa trả, và khoảng thời gian chờ dọn phòng sau khi khách trước rời đi. Trong đó có một trường hợp em đánh giá cao nhất: vẫn phải chặn những phòng đang có khách chưa trả, không phụ thuộc khoảng thời gian. Nó không kiểm tra tính năng em vừa làm, mà kiểm tra rằng tính năng cũ vẫn còn nguyên. Trước đây em nghĩ viết kiểm thử là để chứng minh mã của mình chạy đúng; bây giờ em nghĩ mục đích quan trọng hơn là để người sửa tiếp sau em biết ngay khi họ làm hỏng một hành vi mà họ không biết là có tồn tại.
